\section{Background and Related Work}
\label{sec:related}

\paragraph{Probing and its controls.}
Linear probes track information across network depth \citep{alain2017understanding}, and a probe score reflects both the representation and the learner fitted to it, which motivates matched-capacity probes and control tasks \citep{hewitt2019designing}. The probing literature distinguishes information a readout can recover from information the model uses \citep{belinkov2022probing,ravichander2021probing}; causal erasure tests use by deleting a concept \citep{ravfogel2020null,elazar2021amnesic}, in closed form for every linear reader at once \citep{belrose2023leace}. Our readability grade adopts the control-task idea, and the write test asks a question complementary to erasure: whether the concept's own direction is the one that moves its answer. In medical multimodal models, depth-wise probing maps classification information across towers, connectors, and language layers \citep{zhu2026lost}, and layer-wise probes of MedCLIP expose acquisition and device shortcuts \citep{pedersen2026probes}; we intervene on the decoded clinical direction at the locus where it is read.

\paragraph{Concept directions and steering.}
A concept activation vector is the normal of a classifier separating concept examples from a negative set, and it grades a class by the directional derivative of the logit along that normal \citep{kim2018tcav}; we write the direction forward instead. Such a normal shifts with the layer it is fitted at, entangles with correlated concepts, depends on its negative set, and carries where in the image the concept was read \citep{nicolson2025explaining}. Representation engineering and activation steering write a concept back by changing intermediate representations at inference time \citep{zou2023representation,turner2023steering,panickssery2024caa}, and a probe normal is an appealing intervention because it supplies a label and a signed vector from the same representation. A separability-trained vector can, however, absorb distractor signals and diverge from its concept \citep{pahde2025navigating}: a classifier's weights are a filter rather than the pattern they detect \citep{haufe2014interpretation}, mass-mean directions steer truth judgements better than logistic normals \citep{marks2024geometry}, and the features that identify a concept are largely not those that control the output \citep{arad2025saes}. Benchmarks accordingly score concept detection and steering as separate axes \citep{wu2025axbench}, and alignment search certifies a subspace as causal when it implements an interpretable variable \citep{geiger2024alignments}. These works judge a direction mainly against its own concept. Ownership fixes the estimator and varies the label instead: it asks whether a direction fitted the same way for another concept moves the same answer more. Section~\ref{sec:robust-geometry} runs alternative constructions, a second locus, and token-concentrated writes through this test.

\paragraph{Writing into a vision-language model.}
Image features are linearly decodable in a vision-language model's hidden layers, and targeted edits to them change downstream answers \citep{rajaram2025lineofsight}. Steering one sparse-autoencoder feature of the vision tower and asking the language model what it sees recovers what that feature encodes \citep{ferrando2026explain}; we start from the meaning, fit a direction to a clinical label, and ask whether the write keeps that label's selectivity against the other findings' directions. Causal probing of multimodal models intervenes on the language model's residual stream with difference-in-means vectors and scores each concept against its own effect \citep{deng2026causal}; we write the representation the language model consumes and score each concept against its competitors.

